\section{Limitations and Threats to Validity}
\label{sec:limits}

\emph{Benchmark coverage.} The CIFAR-10 study covers two architectures
(ResNet-18 and VGG-16-BN), two dirty-label attacks (BadNets and Blend), one
target class, and a 5\% poison rate (10\% in the amplification block). It does
not cover deeper or wider networks, vision transformers, ImageNet-scale data, or
clean-label attacks. The cliff could move under a much lower poison rate, where
the residual poison after a 0.9-recall detection is only a handful of samples, or
under a trigger whose effective support is larger than a $4\times4$ patch. We
have no evidence either way, and we do not extrapolate.

\emph{Simulated forget sets.} The primary manipulation constructs the forget set
directly from ground-truth poison labels with requested recall and precision. It
is therefore kinder than a real detector in two ways: false positives are drawn
only from non-target clean data, and the flagged poison is a uniformly random
subset rather than the easiest or hardest cases. The E6 block anchors the grid to
two real detector outputs and reproduces the main conclusion, but two anchors are
not a survey. A detector whose errors are correlated with the poison's
representation, rather than random, could produce different results; our
construction cannot see that correlation.

\emph{Detector coverage.} Only activation clustering and spectral signatures were
swept. Neural Cleanse, ABS, STRIP, and other detectors may have different error
structure. The VGG-16-BN and MobileNetV2 detector runs were degenerate and are
excluded, which means the detector side of the study is, in practice, a
ResNet-18/BadNets study with a cross-check on other cells. The activation
clustering failure at 10\% poison is a property of its thresholds; a detector
with different thresholds might not fail there, although it would then have a
different precision profile.

\emph{Statistical power.} Three seeds per cell give coarse success counts and no
confidence intervals. We report mean and standard deviation and the count of
successful seeds, and we flag the one method (NegGrad+) whose distribution is
bimodal. A larger number of seeds would refine the boundary between 0.90 and
0.95 recall but would not change the qualitative conclusion, which is visible in
every seed. The two-point clean-accuracy budget is a modelling choice; a reader
who prefers one point would move the recipe's $\rstar$ up, and a reader who
prefers three points would move it down. The per-cell numbers are reported so
that the verdict can be recomputed.

\emph{The $r{=}0.95$ clean-accuracy drop.} At $r{=}0.95$ and $p{=}1$, deletion
reaches ASR $2.84\%$ but clean accuracy falls to $84.58\%$, four points below
gold, with no false positives in the forget set. We did not isolate the cause. It
may be an optimisation effect of retraining on a set that still contains 125
poison samples, or a seed-level interaction; the standard deviation across seeds
is small ($0.59$), so it is not seed noise in the usual sense. We report it as
measured and use the joint criterion, which fails the cell regardless of the
cause.

\emph{Unrun blocks.} An adaptive-attacker study (E4) was planned and not
executed. Every result in this paper assumes a static attack that does not know
the defence. An attacker who knows that repair and inoculation will be applied
could try to construct poison that survives relabelling or that is re-activated
by inoculation, and the recipe's tolerance would then be an upper bound in a
stronger sense than we intend. A full hyperparameter sweep for the baselines was
also not run; NegGrad+ and SCRUB were used at their default settings, and a
better-tuned version might perform differently. SISA and fine-pruning were run
under shorter schedules than the frozen protocol, which is why their absolute
clean accuracies are lower; their relative costs are reported with that caveat.

\emph{Radio study.} The trigger is a digital burst applied before normalisation,
not an over-the-air impairment; the classifier is a small one-dimensional CNN;
the per-SNR bins have unequal sample counts; and at SNR below $-10$\,dB the clean
model is near chance. The $\gstar$ values in that region are reported but should
not be interpreted as defence performance. The study is a replication of the
cliff, not a claim about radio deployment.

\emph{Durability probe.} The relearning probe used one hundred fresh poison
samples, a fixed schedule, and no defence-aware optimisation. It was run on the
gold model, not on the recipe's output, so it bounds the best case for erasure
rather than characterising every method. A longer probe or a larger poison budget
would presumably make the backdoor return faster, not slower.

\emph{Theory.} We offer no guarantee. The cliff is an empirical regularity across
the settings we tested, and the paper's claims are correspondingly empirical.
Nothing in the results rules out a removal method, outside our suite, that
provably tolerates a known fraction of missed poison under stated assumptions.
